\section{Conclusion}\label{sec:conclusion}

Synthetic cultural agents compile declared aggregate preference signs into choice policies that can be read on any survey item without naming a group. We evaluated them along the chain from anchor to instrument to respondents, and the evaluation says where they hold and where they fail.

Between anchor and instrument they hold in a limited sense. Training installs the synthetic labels on pairs the adapters never saw, and on trust the adapters keep the anchor's order of countries across a change of scenario bank. A prompt that names the country does not keep that order, and a prompt that states the same signs does, so the declared signs carry the trust order and the evidence cannot credit the trained weights over the prompt. The adapters do not keep the declared dimensions separate. Between instrument and respondents the evidence is weak. A preregistered screen of held-out items, read against a retrospective human criterion, gives a small signal that the registered comparisons do not separate from careful readers of the questionnaire. No dimension passes both the second-bank check and the screen, and the adapters' order of countries does not detectably agree with the respondents' composite. Keeping the anchor's order of countries and approaching respondents' answer distributions are different questions. On a descriptive reading that was not preregistered, the adapters keep the first and lie farther from the survey's answer distributions than the base model, whereas a country-named prompt lies closer with no trust order.

A researcher may use the instrument now as an inspectable, group-level reading of candidate items beside human pretesting, for which it is a complement and no replacement. The test that would license more is prospective: probe-selected items fielded to people after the selection and a human-only comparison rule have both been declared, so that the instrument's selections can be compared with what a researcher would otherwise have fielded. Until that test is run, pre-piloting remains untested.
